A tradução automática baseada em redes neurais tem sido investigada sob diferentes perspectivas, que incluem a construção de modelos, a adaptação a domínios ou línguas com poucos recursos e a avaliação da qualidade das traduções. Os trabalhos discutidos nesta seção foram selecionados por contribuírem para a definição das métricas, dos modelos e dos procedimentos experimentais utilizados neste TCC. Entretanto, eles apresentam objetivos e escopos diferentes, de modo que a comparação entre suas contribuições é necessária para identificar a lacuna abordada neste estudo.

\section{Avaliação automática da qualidade da tradução}

\cite{mukherjee2025detailed} realizam uma análise comparativa entre as métricas neurais BLEURT e BERTScore. O estudo investiga se essas métricas atendem a condições desejáveis para a avaliação automática de tradução, considerando diferentes fenômenos linguísticos e níveis de complexidade textual. Para isso, os autores utilizam os conjuntos TQ-AutoTest, PE\textsuperscript{2}rr e GEC CoNLL, que contemplam, respectivamente, fenômenos linguísticos, erros anotados em nível de palavra e correção gramatical. A análise também contrasta métricas baseadas em representações vetoriais com métricas fundamentadas em correspondência lexical e sintática, como o BLEU.

A principal contribuição desse trabalho está em demonstrar que métricas distintas podem responder de maneira diferente a alterações semânticas, lexicais e gramaticais. Essa conclusão é relevante para o presente TCC, pois justifica a utilização conjunta de BLEU, BERTScore e chrF, em vez da adoção de uma única medida de qualidade. Contudo, o estudo concentra-se na avaliação das próprias métricas e não na comparação sistemática de modelos de tradução inglês--português. Além disso, não tem como objetivo medir o custo computacional dos modelos, nem avaliar a influência de domínios textuais geral e técnico sob um mesmo protocolo experimental.

Em uma perspectiva complementar, \cite{di2025meta} apresentam o corpus MT@BZ, composto por decretos traduzidos por máquina entre italiano e a variedade alemã sul-tirolesa. O corpus contém anotações manuais de erros de tradução no domínio jurídico, produzidas com apoio de especialistas em direito e linguística. A partir dessas anotações, os autores realizam uma meta-avaliação de métricas automáticas, abrangendo métricas baseadas em cadeias de caracteres, como BLEU, chrF e TER, métricas baseadas em embeddings, como BERTScore, e métricas neurais aprendidas, como COMET e BLEURT.

O diferencial de \cite{di2025meta} é a comparação das pontuações automáticas com um padrão de referência construído a partir da gravidade dos erros anotados por especialistas. Assim, o trabalho evidencia a importância de considerar o domínio, o par linguístico e o tipo de erro na interpretação das métricas. Apesar dessa contribuição, seu foco está na validação de métricas em um corpus jurídico específico, e não na análise conjunta da qualidade e da eficiência computacional de modelos de pesos abertos. O trabalho também não compara os modelos MarianMT, M2M-100 e NLLB em textos inglês--português.

Os dois estudos convergem ao mostrar que a avaliação da tradução não deve ser reduzida à sobreposição literal entre a saída e a referência. Entretanto, eles adotam estratégias distintas: \cite{mukherjee2025detailed} enfatizam as propriedades das métricas em diferentes fenômenos linguísticos, enquanto \cite{di2025meta} relacionam as métricas a erros anotados e a julgamentos especializados em um domínio delimitado. Em conjunto, esses trabalhos fornecem fundamentos para a escolha de métricas complementares, mas deixam em aberto como modelos de tradução de arquiteturas e escopos diferentes se comportam simultaneamente em qualidade e consumo de recursos em um cenário controlado.

\section{Modelos, adaptação e eficiência}

\cite{xu2309paradigm} investigam uma estratégia de adaptação de modelos de linguagem grandes à tradução automática. Diferentemente do que a caracterização anterior sugeria, o trabalho não trata da criação de recursos para línguas de Uganda com MarianMT. Os autores propõem o ALMA, construído a partir do LLaMA-2, e aplicam um procedimento de ajuste em duas etapas: inicialmente, utilizam dados monolíngues para ampliar o conhecimento linguístico do modelo e, posteriormente, empregam uma pequena quantidade de dados paralelos de alta qualidade para orientá-lo à tarefa de tradução. Os experimentos utilizam direções dos conjuntos de testes WMT'21 e WMT'22 e avaliam os resultados principalmente por BLEU e COMET.

A contribuição de \cite{xu2309paradigm} está na investigação de uma alternativa aos modelos codificador--decodificador treinados ou ajustados com grandes volumes de dados paralelos. Por outro lado, o estudo avalia uma abordagem de modelo de linguagem grande e concentra-se no treinamento e na qualidade em diferentes direções linguísticas. Não são analisados, no mesmo protocolo, modelos de tradução de pesos abertos especializados e multilíngues, nem são consideradas métricas de uso de CPU, memória RAM e tempo de inferência em execução local.

\cite{moslem2026afrinllb} propõem o AfriNLLB, uma série de modelos leves voltados à tradução de e para línguas africanas. Os modelos são derivados do NLLB-200 de 600 milhões de parâmetros e passam por poda iterativa de camadas, quantização, ajuste fino e destilação de conhecimento. O trabalho reúne e filtra corpora paralelos públicos para 15 pares linguísticos e avalia os modelos com BLEU, chrF++ e métricas COMET especializadas para línguas africanas. Os resultados indicam que a compressão pode reduzir o tempo de inferência mantendo grande parte da qualidade do modelo de referência.

Embora \cite{moslem2026afrinllb} trate diretamente da relação entre qualidade e eficiência, seu objetivo é desenvolver modelos compactos para línguas africanas em ambientes com recursos limitados. Portanto, o estudo não responde à questão de qual modelo de pesos abertos apresenta melhor compromisso entre qualidade e custo computacional na tradução inglês--português sem uma etapa adicional de compressão. Também são diferentes o conjunto de línguas, os corpora, o hardware e o protocolo de avaliação utilizado.

\section{Análise comparativa e lacuna de pesquisa}

Os trabalhos analisados apresentam contribuições importantes, mas abordam partes diferentes do problema. Os estudos de \cite{mukherjee2025detailed} e \cite{di2025meta} concentram-se na confiabilidade e na interpretação das métricas; \cite{xu2309paradigm} investiga uma estratégia de ajuste fino para modelos de linguagem grandes; e \cite{moslem2026afrinllb} enfoca a criação e a compressão de modelos para línguas africanas. Assim, nenhum deles reúne, em um único experimento, a comparação entre modelos de tradução de pesos abertos, a avaliação de diferentes tipos de texto e a medição de recursos computacionais em uma máquina executada apenas em CPU.

A lacuna identificada é, portanto, a ausência de uma comparação controlada entre modelos Transformer com características distintas para a tradução do inglês para o português brasileiro. O presente TCC preenche essa lacuna ao comparar o modelo bilíngue MarianMT, especializado no par inglês--português, com os modelos multilíngues M2M-100 e NLLB. A comparação utiliza os mesmos pares de entrada, as mesmas referências e o mesmo ambiente computacional, contemplando textos gerais da base OpenSubtitles e textos específicos das bases MDN Web Docs e PHP. Dessa forma, o estudo relaciona qualidade, domínio textual e eficiência computacional em um cenário reprodutível de execução local.

A Tabela \ref{tab:trabalhos-correlatos} sintetiza as diferenças entre os trabalhos correlatos e o escopo deste TCC.

\begin{table}[H]
    \centering
    \footnotesize
    \renewcommand{\arraystretch}{1.25}
    \caption{Comparação entre os trabalhos correlatos e o presente TCC}
    \label{tab:trabalhos-correlatos}
    \begin{tabularx}{\textwidth}{>{\raggedright\arraybackslash}p{2.7cm} >{\raggedright\arraybackslash}X >{\raggedright\arraybackslash}X >{\raggedright\arraybackslash}X}
        \toprule
        \textbf{Trabalho} & \textbf{Foco principal} & \textbf{Modelos, dados e métricas} & \textbf{Relação com a lacuna deste TCC} \\
        \midrule
        Mukherjee et al. (2025) & Comparação de métricas neurais e análise de suas limitações. & BLEURT e BERTScore, contrastados com métricas lexicais e sintáticas em três conjuntos de dados. & Não compara modelos de tradução nem mede eficiência computacional em inglês--português. \\
        \addlinespace
        Di Natale et al. (2025) & Meta-avaliação de métricas em traduções jurídicas. & Corpus MT@BZ, italiano--alemão sul-tirolesês, erros anotados por especialistas; BLEU, chrF, TER, BERTScore, COMET e BLEURT. & Relaciona métricas e erros humanos, mas não avalia modelos e domínios diversos sob o mesmo ambiente computacional. \\
        \addlinespace
        Xu et al. (2023) & Adaptação de LLMs para tradução com pouca informação paralela. & ALMA baseado em LLaMA-2; dados monolíngues e paralelos; testes WMT; BLEU e COMET. & Não compara MarianMT, M2M-100 e NLLB nem analisa CPU, RAM e tempo de inferência local. \\
        \addlinespace
        Moslem et al. (2026) & Construção de modelos compactos para línguas africanas. & AfriNLLB derivado do NLLB-200, com poda, quantização e destilação; BLEU, chrF++ e COMET. & Estuda eficiência por compressão, mas não o compromisso entre modelos originais bilíngues e multilíngues em inglês--português. \\
        \addlinespace
        \textbf{Presente TCC} & Comparação de qualidade e eficiência de modelos Transformer de pesos abertos. & MarianMT, M2M-100 e NLLB; textos gerais e específicos em inglês--português; BLEU, BERTScore, chrF, CPU, RAM e tempo. & \textbf{Preenche a lacuna ao combinar comparação de modelos, influência do domínio e custo computacional em um protocolo controlado de execução em CPU.} \\
        \bottomrule
    \end{tabularx}
    \legend{Fonte: elaborado pelo autor.}
\end{table}

Com base nessa comparação, os trabalhos correlatos contribuíram para a seleção de métricas e para a compreensão dos fatores que podem afetar a qualidade da tradução. O presente estudo se diferencia por analisar conjuntamente a qualidade das saídas e o custo de sua produção, além de verificar se o domínio textual altera o desempenho relativo dos modelos. Essa abordagem permite não apenas identificar o modelo com maior pontuação média, mas também discutir a adequação de cada alternativa a cenários em que a qualidade e a disponibilidade de recursos computacionais precisam ser consideradas simultaneamente.
